% !TeX program = pdflatex
\documentclass[11pt,a4paper]{amsart}

\newcommand{\notelanguage}{finnish}
% Shared preamble for course notes and exercise sets.

\usepackage[T1]{fontenc}
\usepackage{lmodern}
\providecommand{\notelanguage}{english}
\usepackage[finnish,english]{babel}
\AtBeginDocument{\expandafter\selectlanguage\expandafter{\notelanguage}}

\newcommand{\notesfinnishlanguage}{finnish}
\ifx\notelanguage\notesfinnishlanguage
\newcommand{\notestheoremname}{Lause}
\newcommand{\noteslemmaname}{Lemma}
\newcommand{\notespropositionname}{Propositio}
\newcommand{\notescorollaryname}{Korollaari}
\newcommand{\notesdefinitionname}{Määritelmä}
\newcommand{\notesexamplename}{Esimerkki}
\newcommand{\notesproblemname}{Tehtävä}
\newcommand{\notesexercisename}{Tehtävä}
\newcommand{\notesremarkname}{Huomautus}
\newcommand{\notessolutionname}{Ratkaisu}
\else
\newcommand{\notestheoremname}{Theorem}
\newcommand{\noteslemmaname}{Lemma}
\newcommand{\notespropositionname}{Proposition}
\newcommand{\notescorollaryname}{Corollary}
\newcommand{\notesdefinitionname}{Definition}
\newcommand{\notesexamplename}{Example}
\newcommand{\notesproblemname}{Problem}
\newcommand{\notesexercisename}{Exercise}
\newcommand{\notesremarkname}{Remark}
\newcommand{\notessolutionname}{Solution}
\fi

\usepackage[a4paper,margin=30mm]{geometry}
\usepackage{microtype}
\usepackage{graphicx}
\usepackage{enumitem}

\newcommand{\exerciseheading}[2]{%
  \par\noindent
  {\large\bfseries #1\par}
  \vspace{0.1em}
  \noindent{\large #2\par}
  \vspace{1.5em}
}

\usepackage{amsmath}
\usepackage{amssymb}
\usepackage{amsthm}
\usepackage{mathtools}
\usepackage{aliascnt}

% Row operations beside the affected row, e.g. \rowop{+2R_1}.
\newcommand{\rowop}[1]{\mathrel{\big|}\,#1}

% Linear systems: the argument gives the number of alignment pairs.
\newenvironment{linearsystem}[1]
{\left\{
  \begin{alignedat}{#1}}
    {
  \end{alignedat}\right.}

  \usepackage[hidelinks,pdfusetitle]{hyperref}

  \numberwithin{equation}{section}
  \setlist{itemsep=0.25em,topsep=0.5em}

  \theoremstyle{plain}
  \newtheorem{theorem}{\notestheoremname}[section]
  \newaliascnt{lemma}{theorem}
  \newtheorem{lemma}[lemma]{\noteslemmaname}
  \aliascntresetthe{lemma}
  \newaliascnt{proposition}{theorem}
  \newtheorem{proposition}[proposition]{\notespropositionname}
  \aliascntresetthe{proposition}
  \newaliascnt{corollary}{theorem}
  \newtheorem{corollary}[corollary]{\notescorollaryname}
  \aliascntresetthe{corollary}

  \theoremstyle{definition}
  \newaliascnt{definition}{theorem}
  \newtheorem{definition}[definition]{\notesdefinitionname}
  \aliascntresetthe{definition}
  \newaliascnt{example}{theorem}
  \newtheorem{example}[example]{\notesexamplename}
  \aliascntresetthe{example}
  \newaliascnt{problem}{theorem}
  \newtheorem{problem}[problem]{\notesproblemname}
  \aliascntresetthe{problem}
  \newtheorem{exercise}{\notesexercisename}

  \theoremstyle{remark}
  \newaliascnt{remark}{theorem}
  \newtheorem{remark}[remark]{\notesremarkname}
  \aliascntresetthe{remark}

  \newenvironment{solution}
  {
  \begin{proof}[\notessolutionname]}
    {
  \end{proof}}

  \usepackage[nameinlink,noabbrev]{cleveref}

  \crefname{theorem}{theorem}{theorems}
  \crefname{lemma}{lemma}{lemmas}
  \crefname{proposition}{proposition}{propositions}
  \crefname{corollary}{corollary}{corollaries}
  \crefname{definition}{definition}{definitions}
  \crefname{example}{example}{examples}
  \crefname{problem}{problem}{problems}
  \crefname{exercise}{exercise}{exercises}
  \crefname{remark}{remark}{remarks}

  \ifx\notelanguage\notesfinnishlanguage
  \crefname{theorem}{lause}{lauseet}
  \crefname{lemma}{lemma}{lemmat}
  \crefname{proposition}{propositio}{propositiot}
  \crefname{corollary}{korollaari}{korollaarit}
  \crefname{definition}{määritelmä}{määritelmät}
  \crefname{example}{esimerkki}{esimerkit}
  \crefname{problem}{tehtävä}{tehtävät}
  \crefname{exercise}{tehtävä}{tehtävät}
  \crefname{remark}{huomautus}{huomautukset}
  \fi

  \newcommand{\NN}{\mathbb{N}}
  \newcommand{\NNone}{\mathbb{N}_1}
  \newcommand{\ZZ}{\mathbb{Z}}
  \newcommand{\QQ}{\mathbb{Q}}
  \newcommand{\RR}{\mathbb{R}}
  \newcommand{\CC}{\mathbb{C}}
  \newcommand{\PP}{\mathbb{P}}

  \DeclareMathOperator{\im}{im}
  \DeclareMathOperator{\rank}{rank}
  \DeclareMathOperator{\syt}{syt}
  \DeclareMathOperator{\pyj}{pyj}
  \DeclareMathOperator{\dom}{dom}
  \DeclarePairedDelimiter{\abs}{\lvert}{\rvert}
  \DeclarePairedDelimiter{\norm}{\lVert}{\rVert}
  \newcommand{\set}[1]{\left\lbrace#1\right\rbrace}
  \newcommand{\maxset}[1]{\max\set{#1}}
  \newcommand{\minset}[1]{\min\set{#1}}
  \newcommand{\powerset}[1]{\mathcal{P}\!\left(#1\right)}


\title{Lukuteoria I\\Harjoitus 2}
\author{}
\date{}

\begin{document}

\exerciseheading{Lukuteoria I}{Harjoitus 2}

\begin{exercise}
  \leavevmode\par
  \begin{enumerate}[label=(\alph*), nosep]
    \item Määritä lukujen \(924\) ja \(630\) alkutekijähajotelmat.
    \item Määritä \(\syt(924, 630)\) sekä alkutekijähajotelmien avulla, että
      Eukleideen algoritmilla.
    \item Esitä \(\syt(924, 630)\) lukujen \(924\) ja \(630\)
      lineaarikombinaationa.
    \item Määritä \(\operatorname{pyj}(924, 630)\).
  \end{enumerate}
\end{exercise}

\begin{solution}
  \leavevmode\par
  \begin{enumerate}[label=(\alph*), nosep]
    \item \(924 = 2^2 \cdot 3 \cdot 7 \cdot 11\), \(630 = 2 \cdot 3^2
      \cdot 5 \cdot 7\).
    \item Alkutekijähajotelman avulla:
      \[
        \begin{aligned}
          \syt(924, 630) &= 2^{\minset{1, 2}} \cdot 3^{\minset{1, 2}}
          \cdot 5^{\minset{0, 1}}
          \cdot 7^{\minset{1}} \cdot 11^{\minset{0, 1}} \\
          &= 2 \cdot 3 \cdot 7 \\
          &= 42.
        \end{aligned}
      \]
      Laskemalla saadaan
      \[
        \begin{aligned}
          924 &= 1 \cdot 630 + 294 \\
          630 &= 2 \cdot 294 + 42 \\
          294 &= 7 \cdot 42 + 0.
        \end{aligned}
      \]
      Täten Eukleideen algoritmin nojalla
      \[
        \syt(924, 630) = \syt(630, 294) = \syt(294, 42) = \syt(42, 0) = 42.
      \]
    \item Kohdan (b) nojalla saadaan
      \[
        \begin{aligned}
          \syt(924, 630) &= 42 \\
          &= 630 - 2 \cdot 294 \\
          &= 630 - 2(924 - 630) \\
          &= 3 \cdot 630 - 2 \cdot 924.
        \end{aligned}
      \]
    \item Alkutekijähajotelmien nojalla
      \[
        \begin{aligned}
          \pyj(924, 630) &= 2^{\maxset{1, 2}} \cdot 3^{\maxset{1, 2}}
          \cdot 5^{\maxset{0, 1}}
          \cdot 7^{\maxset{1}} \cdot 11^{\maxset{0, 1}} \\
          &= 2^2 \cdot 3^2 \cdot 5 \cdot 7 \cdot 11.
        \end{aligned}
      \]
  \end{enumerate}
\end{solution}

\begin{exercise}
  Olkoot \(a\), \(b\) ja \(c\) kokonaislukuja.
  \begin{enumerate}[label=(\alph*), nosep]
    \item Osoita, että jos \(\syt(a, b) = 1\), \(a \mid c\) ja
      \(b \mid c\), niin \(ab \mid c\).
    \item Osoita, että edellisen kohdan tulos ei aina päde, jos
      \(\syt(a, b) > 1\).
  \end{enumerate}
\end{exercise}

\begin{solution}
  \leavevmode\par
  \begin{enumerate}[label=(\alph*), nosep]
    \item Koska \(\syt(a, b) = 1\), on olemassa \(x, y \in \ZZ\) siten, että
      \(1 = ax + by\).
      Koska \(a \mid c\) ja \(b \mid c\), on olemassa \(k, \ell \in \ZZ\)
      siten, että \(c = ka\) ja \(c = \ell b\).
      Saadaan
      \[
        \begin{aligned}
          c &= cax + cby \\
          &= \ell bax + kaby \\
          &= ab(\ell x + ky).
        \end{aligned}
      \]
      Täten \(ab \mid c\).
    \item Olkoon \(a = 2\), \(b = c = 4\). Tällöin \(\syt(a, b) = 2 > 1\),
      \(a \mid c\) ja \(b \mid c\). Kuitenkin \(ab = 8 \nmid 4 = c\).
  \end{enumerate}
\end{solution}

\begin{exercise}
  Ratkaise seuraavat Diofantoksen yhtälöt
  \begin{enumerate}[label=(\alph*), nosep]
    \item \(9x + 21y = 3\)
    \item \(32x - 10y = 5\)
    \item \(-4x + 11y = 7\)
  \end{enumerate}
\end{exercise}

\begin{solution}
  \leavevmode\par
  \begin{enumerate}[label=(\alph*), nosep]
    \item Laskemalla saadaan
      \[
        \begin{aligned}
          21 &= 2 \cdot 9 + 3 \\
          9 &= 3 \cdot 3 + 0.
        \end{aligned}
      \]
      Täten Eukleideen algoritmin nojalla \(\syt(21, 9) = 3\). Diofantoksen
      yhtälöllä on täten ratkaisuja, sillä \(3 \mid 3\). Yksi ratkaisu on
      \((x_0, y_0) = (-2, 1)\). Homogeenisen yhtälön
      \[
        9x + 21y = 0
      \]
      yleinen ratkaisu on
      \[
        (x, y) = (7k, -3k), \quad k \in \ZZ.
      \]
      Täten alkuperäisen Diofantoksen yhtälön yleinen ratkaisu on
      \[
        (x, y) = (-2 + 7k, 1 - 3k), \quad k \in \ZZ.
      \]
    \item Huomataan, että \(\syt(32, 10) = 2 \nmid 5\). Täten Diofantoksen
      yhtälöllä ei ole ratkaisuja.
    \item Laskemalla saadaan
      \[
        \begin{aligned}
          11 &= 2 \cdot 4 + 3 \\
          4 &= 1 \cdot 3 + 1 \\
          3 &= 3 \cdot 1 + 0.
        \end{aligned}
      \]
      Täten \(\syt(11, 4) = 1\). Yhtälöllä on täten ratkaisuja. Ensiksi
      \[
        \begin{aligned}
          1 &= 4 - 1 \cdot 3 \\
          &= 4 - 1(11 - 2 \cdot 4) \\
          &= 3 \cdot 4 - 1 \cdot 11.
        \end{aligned}
      \]
      Kertomalla puolittain luvulla \(7\) saadaan
      \[
        7 = 21 \cdot 4 - 7 \cdot 11 = (-4)(-21) + 11(-7).
      \]
      Yksi ratkaisu on täten \((x_0, y_0) = (-21, -7)\). Homogeenisen yhtälön
      \[
        -4x + 11y = 0
      \]
      yleinen ratkaisu on
      \[
        (x, y) = (11k, 4k), \quad k \in \ZZ.
      \]
      Täten alkuperäisen Diofantoksen yhtälön yleinen ratkaisu on
      \[
        (x, y) = (-21 + 11k, -7 + 4k), \quad k \in \ZZ.
      \]
  \end{enumerate}
\end{solution}

\begin{exercise}
  Olkoot \(a\) ja \(b\) kokonaislukuja sekä \(p\) ja \(q\) erisuuria
  alkulukuja.
  \begin{enumerate}[label=(\alph*), nosep]
    \item Osoita, että jos \(a \equiv b \pmod{p}\) ja
      \(a \equiv b \pmod{q}\), niin \(a \equiv b \pmod{pq}\).
    \item Osoita, että \(p^{q-1} + q^{p-1} \equiv 1 \pmod{pq}\).
  \end{enumerate}
\end{exercise}

\begin{exercise}
  Ratkaise seuraavat kongruenssiyhtälöt
  \begin{enumerate}[label=(\alph*), nosep]
    \item \(3x \equiv 2 \pmod{9}\),
    \item \(15y \equiv 7 \pmod{41}\),
    \item \(30z - 5 \equiv 11 + 2z \pmod{40}\).
  \end{enumerate}
\end{exercise}

\begin{exercise}
  Etsi luvulle \(3^{90} - 1\) ainakin kolme tekijää \(d\), joille pätee
  \(2 < d < 3^{90} - 1\).
\end{exercise}

\end{document}
